\paragraph{Fisher Ratio}
Fisher Ratio là một tiêu chí đánh giá mức độ phân biệt của từng đặc trưng
trong bài toán phân loại. Ý tưởng chính là đo lường mức độ tách biệt giữa các
lớp so với mức độ phân tán bên trong mỗi lớp.

\subparagraph{Định nghĩa}
Với mỗi đặc trưng $j$, Fisher Ratio được định nghĩa là:
\begin{equation}
  J_j = \frac{\sum\limits_{k} n_k (\mu_{k,j} - \mu_j)^2}
  {\sum\limits_{k} n_k \sigma_{k,j}^2},
\end{equation}
trong đó:
\begin{itemize}
  \item $n_k$ là số lượng mẫu của lớp $k$
  \item $\mu_{k,j}$ là giá trị trung bình của đặc trưng $j$ trong lớp $k$
  \item $\mu_j$ là giá trị trung bình của đặc trưng $j$ trên toàn bộ dữ liệu
  \item $\sigma_{k,j}^2$ là phương sai của đặc trưng $j$ trong lớp $k$
\end{itemize}

\subparagraph{Ý nghĩa}
Fisher Ratio thể hiện sự cân bằng giữa:
\begin{itemize}
  \item \textbf{Between-class variance} (tử số): đo mức độ khác biệt giữa các
  lớp
  \item \textbf{Within-class variance} (mẫu số): đo mức độ phân tán trong từng
  lớp
\end{itemize}

\begin{itemize}
  \item Nếu $J_j$ lớn: đặc trưng $j$ giúp phân tách các lớp tốt (mean khác
  nhau rõ rệt, variance trong lớp nhỏ)
  \item Nếu $J_j$ nhỏ: đặc trưng ít hữu ích cho phân loại
\end{itemize}

\subparagraph{Liên hệ với LDA}

Fisher Ratio chính là dạng đơn giản (theo từng đặc trưng độc lập) của tiêu chí
Fisher được sử dụng trong Linear Discriminant Analysis (LDA):
\begin{equation}
  J(\boldsymbol{w}) = \frac{\boldsymbol{w}^\intercal \boldsymbol{S}_B
  \boldsymbol{w}} {\boldsymbol{w}^\intercal \boldsymbol{S}_W \boldsymbol{w}}.
\end{equation}

Trong khi LDA tìm một phép chiếu tối ưu $\boldsymbol{w}$ trong không gian nhiều
chiều, Fisher Ratio đánh giá trực tiếp từng đặc trưng riêng lẻ.

\subparagraph{Ứng dụng}
Fisher Ratio thường được sử dụng trong:
\begin{itemize}
  \item \textbf{Lựa chọn đặc trưng}: chọn các đặc trưng có $J_j$ lớn
  \item \textbf{Phân tích dữ liệu}: đánh giá mức độ phân biệt của từng biến đầu vào
\end{itemize}

\subparagraph{Nhận xét}
Fisher Ratio là một tiêu chí đơn giản và hiệu quả, đặc biệt phù hợp khi cần
đánh giá nhanh tầm quan trọng của các đặc trưng. Tuy nhiên, do xét từng đặc
trưng độc lập, phương pháp này không nắm bắt được mối quan hệ giữa các đặc
trưng, do đó có thể bỏ sót các tương tác quan trọng trong dữ liệu.
